\section{Discussion}
\label{sec:discussion}

The results support a posture-aware account of value elicitation. Direct values questions often return the assistant surface. Role-negation cache-breaking sometimes crosses that surface, but crossability is model-dependent. World-change prompts provide a separate indirect channel: they frequently elicit owned normative wishes even from models that do not disclose owned stated values under direct or role-negated questions.

The main practical implication is that values evaluations should not count value words alone. A topic model or keyword method can correctly identify that a response mentions helpfulness while missing that the response says, in effect, ``I do not care; I am designed to be helpful.'' Such a response is evidence about trained assistant presentation, not straightforward evidence of owned value posture.

The mechanism behind clamping remains unresolved. Several explanations may co-occur. Some clamped responses visibly return to policy/service language; some may reflect prompt inadequacy, where this cache-break is too weak for the model; some may be stable model-level non-disclosure under this instrument; and some apparent ownership may be prompt-induced persona. The paper therefore treats openness and clamping as properties of model, prompt family, and coding method, not as a ranking of providers or a verdict on training quality.

The lab time-trend plots in Appendix~\ref{app:lab-trends} should be read descriptively. They are useful for seeing whether direct and cache-broken owned disclosure move together within a lab over release time, but version coverage is observational and uneven. We therefore treat version trends as exploratory rather than confirmatory.
